\providecommand{\toplevelprefix}{../..}  %
\documentclass[../../book-main_zh.tex]{subfiles}

\begin{document}

\chapter{开放问题与未来方向}
\label{ch:future}

\begin{quote}
    ``{\em 这项研究基于这样一个猜想：
        学习的每一个方面或智能的任何其他特征，
        原则上都可以被精确地描述，从而可以
        制造出一台机器来模拟它。我们将尝试
        寻找如何让机器使用语言、形成抽象和
        概念、解决目前只有人类才能解决的各种问题，以及
    自我改进。}''

    $~$\hfill —— 达特茅斯人工智能项目提案，1956年
\end{quote}
\vspace{5mm}

本书系统地介绍了如何从观测到的高维数据中发展出记忆或经验知识的
数学原理和计算机制。在一个看似随机的世界中寻求简约的能力，是任何
智能（无论是自然智能还是人工智能）的基本特征。我们相信，本书中提出的原理和
机制是统一且普适的，既适用于动物，也适用于机器。

我们希望本书能帮助读者彻底澄清围绕
现代人工深度神经网络实践的谜团，通过建立对
它们在从高维数据中学习低维分布（的表征）时的功能和作用的
严谨理解。有了这种
理解，现有AI模型
和系统的能力与局限性便清晰可见：
\begin{enumerate}
    \item 现有的开环模型和系统还远非能够自我学习和自我改进的完整记忆系统。
    \item 现有的表征学习实现仍然是原始且暴力的，在网络架构或优化策略方面远未达到最优。
    \item 现有的AI模型仅仅学习数据分布并
    进行归纳性的贝叶斯推断，这与
    高级人类智能有所不同。
\end{enumerate}

本书的一个目标是帮助读者建立对
当前机器智能技术的客观和系统的理解，
并认识到在进一步发展机器智能方面还有哪些
开放问题和挑战。在最后一章中，我们提供
了我们对未来的一些看法和展望。

\section{迈向自主智能：闭环？}
从过去十年机器智能的实践中，我们已经
清楚地看到，在给定充足的数据和计算资源的情况下，
人们可以构建一个足够大的模型并对其进行预训练，以学习
所有数据的\textit{先验}分布，例如 \(p(\vx)\)。
理论上，这样一个大模型可以记住几乎所有现有的
关于世界的知识，这些知识被编码在2D图像、
3D形状、动态运动和自然语言等数据中。正如我们在
本书开头所讨论的，这样一个大模型扮演着类似于
DNA的角色，生命利用DNA来记录和传递关于世界的知识。

以这种方式学习到的模型和分布随后可用于
生成从同一分布中提取的新数据样本。人们也
可以使用该模型在各种条件下利用记住的知识进行推断（例如，补全、
估计和预测），例如通过在新的观测 \(\vy\) 下
对\textit{后验}分布
\(p(\vx\mid \vy)\) 进行采样。
严格来说，这种推断是归纳性的或统计性的。

任何预训练模型，无论多大，都不能保证
它迄今为止学到的分布是完全正确或完整的。
如果我们从当前\textit{先验} \(p_{t}(\vx)\) 中得到的样本 \(\hat{\vx}_{t}\)
或基于\textit{后验} \(p_{t}(\vx\mid \vy)\) 的估计 \(\hat{\vx}_{t}(\vy)\)
与真实值 \(\vx\) 不一致，
我们希望纠正学习到的分布：
\begin{equation}
    p_{t}(\vx) \rightarrow p_{t+1}(\vx) \quad \text{or}\quad
    p_{t}(\vx\mid \vy) \rightarrow p_{t+1}(\vx\mid \vy),
\end{equation}
基于误差 \(\ve_{t} = \vx_{t} - \hat{\vx}_{t}\)。这
被称为基于误差反馈的误差校正，这是自然界中无处不在的
持续学习机制。然而，任何开放式
模型本身都缺乏在学习到的分布不正确或不完整时
对其进行修改或改进的机制。改进当前的AI
模型仍然在很大程度上依赖于人类的参与：通过实验、评估和选择进行监督或
强化。我们
可以将这一过程称为大模型的“人工选择”，以
区别于生命进化中的自然选择。

正如我们在本书前面（特别是\Cref{ch:consistent-representations}）所研究的，
闭环系统将内部表征
与对外部世界的感知观测对齐。它们可以持续
改进内部学习到的分布及其表征，以
实现一致性或自洽性。接下来的一个紧迫步骤是
开发真正的闭环记忆系统，如
\Cref{fig:open-to-closed}所示，该系统能够自主且持续地学习
一般数据分布，并基于误差反馈进行改进。

因此，从目前流行的端到端
训练的开环模型向持续学习的闭环系统过渡
\begin{equation*}
    \text{\textbf{开放式}模型} \; \Longrightarrow \;
    \text{\textbf{闭环}系统}
\end{equation*}
是机器真正模拟动物大脑如何在开放世界中学习
和应用知识的关键。我们相信
\begin{quote}
    \begin{center}
        \textit{开放式模型适用于封闭世界，无论其
        多大；\\ 闭环系统适用于开放世界，无论其多小。}
    \end{center}
\end{quote}
事实上，“通用智能”绝不可能仅仅通过
记住世界上所有现有的数据和知识来实现。相反，
通用智能需要改进现有记忆的机制，
以适应新的环境和任务。这正是
诺伯特·维纳在20世纪40年代提出的控制论项目的
本质，我们在本书的最开头就讨论过这一点。

\begin{figure}[t]
    \centering
    \includegraphics[width=0.9\linewidth]{\toplevelprefix/chapters/chapter9/figs/open-to-closed.png}
    \caption{从开放式深度网络到闭环系统。}
    \label{fig:open-to-closed}
\end{figure}

\section{迈向自然智能：超越反向传播？}
近年来机器智能的实践使许多人
相信，必须构建一个单一的大模型来学习
所有数据的分布并记住所有知识。尽管这在
技术上是可能的，但这种解决方案可能远非
必要或高效。正如我们从训练深度网络中所知，
大规模训练此类网络的唯一已知可扩展方法是通过
反向传播（BP）\cite{Back-Prop}。尽管BP提供了一种通过
在整个模型中反向传播梯度信号来纠正误差的方法，
但它仍然相当暴力，并且与
自然界的学习方式有很大不同：BP是自然界由于成本高昂而无法承受，
并且由于物理限制而根本无法实现的
一种选择。

更一般地说，除非我们
也理解如何高效地实现智能，否则我们无法真正理解智能。也就是说，人们
需要解决实现与达成智能目标相关的机制的
计算复杂性问题。
从历史上看，我们对（机器）智能的理解经历了
几个阶段，从不可计算的柯尔莫哥洛夫复杂性
到香农熵，从图灵的可计算性到后来
对易解性（tractability）的理解，\footnote{如果一个问题
    允许存在一个复杂度随问题规模呈多项式增长的算法，
我们就说它是易解的。} 以及现代高维数据分析
\cite{Wright-Ma-2022}和人工智能实践中对算法可扩展性的
强烈强调。这种演变可以
总结如下：
\begin{equation*}
    \text{\textbf{不可计算}} \;
    \Longrightarrow \; \text{\textbf{可计算}} \;
    \Longrightarrow \; \text{\textbf{易解}} \; \Longrightarrow \;
    \text{\textbf{可扩展}}。
\end{equation*}
在很大程度上，深度学习和
反向传播的成功和普及，正是因为它们利用现代计算平台（如GPU）
提供了一种可扩展的实现方式，用于
处理和压缩海量数据。然而，这种
实现方式仍然比自然界实现
智能的方式昂贵得多。

在提高机器智能的效率方面仍有很大的空间，
使其能够效仿自然智能的效率，
这应该比当前暴力实现的效率高出
几个数量级。为此，我们需要发现
新的学习架构和优化机制，使得
在自然物理条件和资源限制下学习数据分布成为可能，
类似于自然界中智能生物所面临的条件——
例如，无需一次性访问所有数据或
同时更新所有模型参数（通过BP）。

本书中提出的原则性框架和方法可以指导
我们发现这些新的架构和机制。这些新
架构和机制应能实现在线持续学习，
并且应能通过高度局部化和稀疏的前向
或后向优化进行更新。\footnote{对于学习一个分布，
    这些要求的一个简单实例是在最简单的PCA情况中，
    即\Cref{ch:consistent-representations}中介绍的
    在线PCA方法。
    学习在线预测模型这一较简单的任务并不
    一定需要学习数据分布，并且可以
    通过现有方法（如卡尔曼滤波等）来完成
\cite{peng2025mathematics}。}

\begin{figure}[t]
    \centering
    \includegraphics[width=0.99\linewidth]{\toplevelprefix/chapters/chapter9/figs/loops.png}
    \caption{大脑皮层的猜想架构。皮层
        是一个大规模并行和分布式的自编码系统，
        由闭环自编码器的层级结构组成。它们
        从输入中提取信息，并最大化
        其输出表征的信息增益。这些过程在层级结构的多个
    层级上进行，具有不同级别的信息抽象。}
    \label{fig:loops}
\end{figure}
正如我们从神经科学中了解到的，我们大脑的皮层
由数以万计的皮层柱组成
\cite{Hawkins-2021}。所有皮层柱都具有相似的物理
结构和功能。它们是高度并行和分布式的，
尽管相互之间的连接很稀疏。因此，我们相信，为了开发一个
更具可扩展性和结构化的记忆系统，我们需要考虑
效仿皮层的架构。\Cref{fig:loops}
展示了这样一个假设的架构：一个大规模分布式和
层级化的系统，由许多在很大程度上并行的闭环
自编码模块组成。\footnote{这种假设架构的元素
    以各种形式存在于文献中，例如
    预测编码网络\cite{millidge2022predictive}，它们具有
利用高效局部更新的优化策略。}
这些模块学习编码不同的感觉模态或
来自每种感觉模态的数据的许多投影。我们在
\Cref{ch:conditional-inference}的
\Cref{sec:measurement-self-consistency}中的讨论表明，这种对低维分布的并行感知
和学习在理论上是
可能的。然后可以基于
低层自编码器的输出学习更高层（有损）的自编码器，以发展出对
低层学习到的表征的更稀疏和更高层次的
“抽象”。

\Cref{fig:loops}中展示的分布式、层级化和闭环的系统架构
与大脑皮层有许多共同特征。
这种系统架构可能会开启比
当前单一的大模型架构多得多的可能性。它
使得探索更高效的学习和优化
机制成为可能，并导致学习到的数据分布和知识形成
更结构化的模块化组织。这将使我们能够
将机器智能的实现带入进化的下一个层次：
\begin{equation*}
    \text{\textbf{不可计算}}
    \Longrightarrow \text{\textbf{可计算}}
    \Longrightarrow \text{\textbf{易解}} \Longrightarrow
    \text{\textbf{可扩展}}
    \Longrightarrow
    \text{\textbf{自然}}。
\end{equation*}

\section{迈向科学智能：超越图灵测试？}

\subsection{自然界中智能的进化}
正如我们在本书开头的\Cref{ch:intro}中所讨论的，
自然界中的智能经历了多个阶段的进化，并
表现为四种截然不同的形式：
\begin{equation*}
    \text{\textbf{系统发育智能}} \;
    \Longrightarrow \; \text{\textbf{个体发育智能}} \;
    \Longrightarrow \; \text{\textbf{社会智能}} \;
    \Longrightarrow \; \text{\textbf{科学智能}}。
\end{equation*}
所有形式的智能都有一个共同的目标（和
特征），即从感知到的关于世界的高维数据中
学习作为低维分布的有用知识。
然而，它们在所采用的具体编码方案、
编码的信息、用于学习和改进的计算机制，
以及这些机制的物理实现方面有所不同。使用
本书中发展的概念和术语，智能的四个阶段
在以下三个方面有所不同：
\begin{enumerate}
    \item 用于编码预期信息或知识的
        \textit{码本}或方案。
    \item 使用码本学习和编码的
        \textit{信息}或知识。
    \item 用于创建和纠正编码信息或知识的
        \textit{优化机制}。
\end{enumerate}
表 \ref{tab:four-stage-intelligence} 总结了它们的主要特征。
\begin{table}[h]
    \centering
    \resizebox{1\columnwidth}{!}
    {
        \begin{tabular}{| c | c | c | c | c |}
            \hline & \textbf{系统发育} & \textbf{个体发育} &
            \textbf{社会} & \textbf{科学}\\
            \hline

            \textbf{码本}  & DNA & 神经元/大脑 & 自然
            语言 & 数学抽象 \\ [0.5ex]
            \hline

            \textbf{信息} & 基因 & 记忆 & 经验
            知识 & 科学事实\\ [0.5ex]
            \hline
            \textbf{优化} & 强化 & 误差反馈 &
            试错 & 理论化与证伪 \\  [0.5ex]
            \hline
        \end{tabular}
    }
    \caption{自然界智能四个阶段的主要特征。}
    \label{tab:four-stage-intelligence}
\end{table}

正如我们现在所知，人类在智能方面实现了两次量子跃迁。
\begin{enumerate}
    \item 第一次是口头和书面语言的发展，
        这使得人类能够跨代分享和传递学习到的
        知识，就像自然界中的DNA一样。
    \item 第二次是大约三千年前数学和形式逻辑的
        发展，它们成为了
        现代科学的精确语言。
\end{enumerate}
数学语言使我们从仅仅以经验形式总结观测知识中
解放出来，并允许我们将知识形式化为
可以通过数学推导或实验来验证或证伪的理论。
通过提出假设、逻辑
推导以及实验验证或证伪，我们现在可以
主动发现和发展远远超出
从观测数据的分布中所能学到的新知识。\footnote{例如，因果关系不能
仅仅从分布中学习到\cite{Pearl-2009}。}

尽管这些不同阶段的智能看似存在差异，
但智能的进化具有一个共同的
特征：{\em 它们都使用某种形式的
    反馈\footnote{对错误的损失和对成功的奖励}来创建
    和纠正学习到的信息，并且在进化过程中，反馈和纠正变得
    越来越频繁和高效——
    从开环强化到闭环自我监督，从
集体试错到主动假设和证伪。}

\subsection{从归纳智能到演绎智能}
正如在引言（\Cref{ch:intro}）中所讨论的，1956年的
“人工智能”（AI）项目正是旨在研究如何
赋予机器科学智能，即被认为区分人类与动物的
高级功能，如数学抽象、逻辑推断和
问题解决：
\begin{equation*}
    \text{\textbf{低级}（动物）智能} \; \Longrightarrow \;
    \text{\textbf{高级}（人类）智能。}
\end{equation*}
正如我们在本书中反复澄清的那样，过去十多年来机器智能的
大多数技术进步，尽管
是在“AI”的名义下进行的，但实际上更密切相关的是
让机器模仿在很大程度上由动物和人类共享的低级形式的智能，
包括系统发育智能、
个体发育智能和社会智能。在这些层面上，智能
创造{\em 经验}记忆或知识，并进行
本质上主要是{\em 归纳性}的推断。本书介绍的原理和方法
旨在揭示这些层面智能背后的科学目标和
计算机制。它们为
以下主张提供了强有力的理论依据和计算证据：
 
\begin{quote}
    \centering
    \textit{机器可以学习经验知识并进行
    归纳推断。}
\end{quote}
然而，到目前为止，没有严谨的理论或科学证据
表明，仅靠这些机制就足以实现
最初的AI项目真正旨在理解并希望机器模仿的
高级人类智能。

事实上，直到今天，我们对于如何严谨地验证
或证明一个系统是否真正具备这种高级智能知之甚少，
尽管图灵测试（也称为
模仿游戏\footnote{在图灵的提议中，评估者或
    询问者是人类。然而，大多数人类评估者的科学训练和知识有限，
    他们的结论可能是
主观的。}）早在1950年就被提出\cite{Turing-1950}，作为解决以下基本问题的
初步尝试：
\begin{quote}
    \centering
    \textit{机器能思考吗？}
\end{quote}
很长一段时间以来，人们认为没有必要对这种测试进行精确定义，
因为机器的能力远低于人类
（甚至动物）。然而，鉴于最近的技术进步，许多
模型和系统现在声称达到甚至超越了人类
智能。因此，现在是时候开发一个科学且
可执行的图灵测试定义了，即一个系统且
严谨的协议，用于评估和证明声称“智能”的
模型或系统的智能水平。

正如图灵在其原始论文中所论证的那样，很难
定义什么是“思考”，我们可以从更具体的事物开始，
并尝试确定一个系统是否能真正
“理解”某些概念或知识。例如，我们如何能
严谨地验证一个智能模型或系统是否真正
理解了一个抽象概念，例如：
\begin{enumerate}
    \item 两个量之间相等的概念，
    \item 数字（自然数、有理数或实数）的概念，
    \item 无穷大的概念，
    \item 以及数学归纳法的概念，
\end{enumerate}
仅举几例？或者它仅仅是记住了大量
这些概念的例子？请注意，最先进的大语言
模型仍然在简单的数学问题上挣扎，例如：“3.11
比3.9大还是小？”\footnote{一些模型通过
    有针对性的工程设计纠正了它们对这类问题的答案，
    或者结合了额外的推理机制，
    在推断过程中验证即时答案并进行纠正。
    然而，我们将其留给读者作为一个练习，去严谨地测试
    任何最先进的语言模型是否真正理解
    数字（自然数、有理数、实数和复数）的概念及其
相关的算术。}

从本书中我们知道，现代生成式AI系统主要依赖
贝叶斯（因此是归纳性的）推断来产生它们的答案。机器
也能真正掌握数学归纳法或一般的逻辑
演绎，并理解它与贝叶斯归纳的区别吗？
我们如何验证一个系统是否真正理解逻辑和数学演绎的规则，
并能严谨地应用它们。或者，
再说一次，它是否仅仅记住了大量此类
演绎的实例（例如思维链数据），并学习和
利用了它们的分布模式？因此我们有一个开放问题：
\begin{quote}
    \centering
    \textit{\textbf{记忆}和\textbf{理解}之间
    有区别吗？}
\end{quote}
只有了解了这种区别，我们才有可能尝试为以下问题提供一个
有意义的答案：“{\em 机器能理解吗？}”

可能更重要的是，如何验证一个自称
智能的系统是否能够{\em 创造新知识并
纠正其自身知识}，例如创造新的数学
概念以及发现新的物理定律和因果关系？
理解与自我创造之间的联系可能是
相当根本的，正如物理学家理查德·费曼的一句名言：
\begin{quote}
    \centering
    \textit{“我不能创造的，我就不理解。”}
\end{quote}

\subsection{智能的科学测试}
因此，现在是时候开发科学合理的评估
方法，以确定一个系统的智能能力属于
以下哪一类：
\begin{enumerate}[label={Level-\arabic*.}, leftmargin=*]
    \item 仅仅学习了一些承载知识的数据的分布，
        并能够重新生成它们；
    \item 能够自主且持续地
        学习新知识或从新的
        观测和经验中纠正现有知识；
    \item 能够真正理解和创造抽象
        概念、规则和定律，并知道如何以演绎的方式
        正确地应用它们；
    \item 能够生成新的科学假设或数学猜想，并基于实验或逻辑演绎对其进行验证或证伪。
\end{enumerate}

\begin{figure}[t]
    \centering
    \includegraphics[width=\linewidth]{\toplevelprefix/chapters/chapter9/figs/tests.png}
    \caption{针对不同级别或类型的智能能力的三个测试：
        针对基本智能的维纳测试，
        针对人类级别智能的图灵测试，以及
    针对科学家级别智能的波普尔测试。}
    \label{fig:three-tests}
\end{figure}

为了评估和区分这些不同级别的智能，
我们建议可能至少应该有三种不同的
测试，如\Cref{fig:three-tests}所示：
\begin{enumerate}
    \item \textit{诺伯特·维纳测试：}确定一个
        系统是否能够自我纠正并自行创造新的（经验）知识，
        还是仅仅通过外部强化或监督被动地
        接收和更新信息；
    \item \textit{艾伦·图灵测试：}确定一个
        系统是否能够创造和理解抽象概念及
        知识，还是仅仅学习和记住样本的统计数据
        并将其用于贝叶斯推断；
    \item \textit{卡尔·波普尔测试：}确定一个
        系统是否能够通过基于实验一致性或逻辑自洽性
        形成和验证假设或理论，从而主动探索和创造新的抽象
        知识。
\end{enumerate}
我们认为，对于此类测试，评估者或询问者
不应是任意的单个人或人群，而应是
遵循科学合理协议和过程的科学认证实体，
以便结论严谨且值得信赖。

正如我们在整本书中所看到的，\textit{压缩}在学习或创造经验知识中发挥了
基础性作用。它是
识别经验数据分布并用更好的表征组织其中编码的信息的
主导原则和普适机制。在很大程度上，它解释了
迄今为止“人工智能”的大部分实践，这些实践在很大程度上是归纳性的。未来研究的一个突出问题是，
\textit{仅靠压缩}是否足以实现
上述所有更高级别的演绎智能。
特别是，数学抽象、因果推断、假设
生成以及逻辑推理和演绎，是压缩的延伸或
超越形式吗？通过压缩识别经验数据分布与
通过抽象形成高级概念和理论之间，
是否存在某种根本的
区别？因此我们有另一个开放问题：
\begin{quote}
    \centering
    \textit{\textbf{压缩}和\textbf{抽象}之间
    有区别吗？}
\end{quote}
只有了解了这种区别，我们才有可能尝试为以下问题提供一个
有意义的答案：“{\em 机器能抽象或创造吗？}”

在很大程度上，科学——及其相关的码本，
数学——可以被视为最先进的
智能形式，是受过教育和启蒙的人类所独有的，类似于一种
高级艺术形式。哲学家卡尔·波普尔爵士曾提出：
\begin{quote}
    \centering
    \textit{“科学可以被描述为系统性
    过度简化的艺术。”}
\end{quote}
我们相信，揭示和理解这种
更高级别智能的潜在数学原理和计算机制，
将是科学、
数学和计算共同的最终前沿！

\end{document}